
The \emph{Markovian noise covariance matrix} captures both the marginal variance and the temporal correlations of the noise:
\begin{equation}
\Sigma_{\varepsilon}^{(\mathrm{M})}
  = \mathbb{E}_{\pi} [\varepsilon(Z_{0}) \varepsilon(Z_{0})^{\mathsf{T}}]
    + \sum_{\ell=1}^{\infty} \left( \mathbb{E}_{\pi} [\varepsilon(Z_{0}) \varepsilon(Z_{\ell})^{\mathsf{T}}]
        + \mathbb{E}_{\pi} [\varepsilon(Z_{\ell}) \varepsilon(Z_{0})^{\mathsf{T}}] \right).
\label{eq:noise-cov}
\end{equation}
The series is absolutely convergent under Assumptions 1 and 3. Indeed, for
centered bounded $\varepsilon$,
$\lVert \mathsf{Q}^{\ell} \varepsilon \rVert_{\infty} \le
2 \lVert \varepsilon \rVert_{\infty} (1 / 4)^{\left\lfloor \ell / t_{\mathrm{mix}} \right\rfloor}$, and hence
\[
\lVert \mathbb{E}_{\pi} [\varepsilon(Z_{0}) \varepsilon(Z_{\ell})^{\mathsf{T}}] \rVert
  \le \lVert \varepsilon \rVert_{\infty} \lVert \mathsf{Q}^{\ell} \varepsilon \rVert_{\infty}
  \le 2 \lVert \varepsilon \rVert_{\infty}^{2} (1 / 4)^{\left\lfloor \ell / t_{\mathrm{mix}} \right\rfloor}.
\]
The same bound applies to the transposed covariance term, so the covariance
series converges absolutely in operator norm.
This matrix is the limiting covariance in the Markov chain CLT for the partial sums $n^{-1/2} \sum_{t=0}^{n-1} \varepsilon(Z_{t})$ (cf. Douc et al., 2018, Theorem 21.2.10).

The \emph{asymptotically optimal covariance matrix} is given by
\begin{equation}
\Sigma_{\infty}
  = \overline{A}^{-1} \Sigma_{\varepsilon}^{(\mathrm{M})} \overline{A}^{-\mathsf{T}}.
\label{eq:asymp-cov}
\end{equation}
This is the covariance target attained by the averaged linearized recursion. We call it optimal in the usual averaged-SA sense; a full Hájek--Le Cam optimality statement would require an additional local-asymptotic experiment argument, which is not part of this thesis.

The \emph{Lyapunov equation} plays a central role in the contraction analysis. For any $P = P^{\mathsf{T}} \succ 0$, there exists a unique $Q = Q^{\mathsf{T}} \succ 0$ satisfying $\overline{A}^{\mathsf{T}} Q + Q \overline{A} = P$. Defining $a = \lambda_{\min} (P) / (2 \lVert Q \rVert)$ and $\kappa_{Q} = \lambda_{\max} (Q) / \lambda_{\min} (Q)$, the key contraction property holds: for all $\alpha \in [0, \alpha_{\infty}]$,
\begin{equation}
\lVert I - \alpha \overline{A} \rVert_{Q}^{2} \le 1 - \alpha a.
\label{eq:contraction}
\end{equation}

Since the iterates $\theta_{k}^{(\alpha)}$ alone are generally not Markovian (due to the Markovian noise), we consider the \emph{joint process} $(\theta_{k}^{(\alpha)}, Z_{k+1})$ with kernel
\[
\overline{\mathsf{P}}_{\alpha} f(\theta, z)
  = \int_{\mathsf{Z}} \mathsf{Q}(z, d z') f(F_{z}(\theta), z'),
\]
where $F_{z} (\theta) = (I - \alpha A(z)) \theta + \alpha b(z)$. Thus the current
second coordinate $z$ is the observation used to update $\theta$, and the next
coordinate $z'$ is carried forward to the following step. Under Assumptions
1--3, this joint chain admits a unique invariant distribution $\Pi_{\alpha}$ for
sufficiently small $\alpha > 0$ (Levin et al., 2025).
